\documentclass[10pt,a4paper]{amsart}
\usepackage[margin=19mm]{geometry}
\usepackage{amsmath,amssymb,mathtools}
\usepackage[scr=rsfs,cal=euler]{mathalfa}
\usepackage{xfrac}
\usepackage[unicode,hidelinks,bookmarks=false]{hyperref}
\newcommand{\ie}{i.e.\ }
\newcommand{\cf}{cf.\ }
\newcommand{\resp}{respectively\ }
\newcommand{\Subsection}{\subsection}
\providecommand{\nobreakdash}{\nobreak\hbox{-}}
\setlength{\emergencystretch}{2em}
\multlinegap0pt
\usepackage{bm}
\usepackage{bbm}
\usepackage{dsfont}
\usepackage{xspace}
\usepackage{cancel}
\usepackage{mathtools}
\usepackage{amsthm}
\makeatletter
\@ifundefined{theorem}{\newtheorem{theorem}{Theorem}}{}
\@ifundefined{lemma}{\newtheorem{lemma}[theorem]{Lemma}}{}
\makeatother
\newcommand{\bF}{\mathbb{F}}
\title[Periodicity of traces of Hecke operators modulo prime powers]{Periodicity of traces of Hecke operators modulo prime powers}
\author{Jonas Bergstr\"om, Sjoerd de Vries}
\date{Source: arXiv:2601.03029v2 (CC BY 4.0)}
\begin{document}
\maketitle

\noindent\emph{Source and licence.} Periodicity of traces of Hecke operators modulo prime powers. Version arXiv:2601.03029v2 (CC BY 4.0), distributed under the Creative Commons Attribution 4.0 International licence, as recorded in the arXiv licence metadata for this exact version and shown on the abstract page https://arxiv.org/abs/2601.03029v2. Full rights notice: licenses/arxiv-2601.03029-CC-BY-4.0.txt.\par\smallskip

\noindent\textit{Editorial scope.} Counted unit: the Lemma whose source label is \texttt{lem:denom\_divisible} in the article (source0.tex lines 634--637, source label \texttt{lem:denom\_divisible}), delivered together with the article's own complete original proof of it (source0.tex lines 639--670, one \texttt{proof} environment of 32 lines). No mathematics is written, repaired or re-derived: statement and proof are the article's own text line for line. Required context retained uncounted: lem:rootofunity (lines 548--553); lem:div\_upgrade (lines 611--613). Every definition, display and label of those lines is retained unchanged. Omitted from this package and not redistributed: 1--547, 554--610, 614--633, 638--638, 671--1275. Nothing inside a retained excerpt is omitted or altered: every retained body is delivered exactly as the article's lines are written, so no omission receipt of any kind accompanies this package. Citations: the retained excerpts carry no citation command at all, so no bibliography entry of the article is transcribed and no reference is redistributed. Reference adaptation: none was needed and none was made; every reference inside the delivered unit resolves inside it, so no unresolved reference or citation is produced. Printed slips kept literal: no printed word, symbol, hypothesis or number of the counted statement or of its proof was altered. Layout and class: the article's own class is \texttt{article} with packages including inputenc, url, hyperref, geometry, babel, biblatex, mathtools, amsmath, amssymb, mathrsfs, dsfont, graphicx, booktabs, array; the wrapper is the lane's pinned \texttt{amsart} editorial wrapper at 10pt on A4, which restates the theorem-like environments the retained text needs and adds the article's own macro definitions that the retained text needs (\texttt{\textbackslash bF}), with extra wrapper packages (bm, bbm, dsfont, xspace, cancel, mathtools). The delivered numbering is the unit's own. Assets: no retained span contains a figure environment, a graphics-inclusion command, a PDF-page inclusion, a table or a TikZ picture; no image file is copied into this package, the package declares no asset dependency and no image file is redistributed with it. Licence: version arXiv:2601.03029v2 of this article is distributed under the Creative Commons Attribution 4.0 International licence, as recorded in the arXiv licence metadata for this exact version and shown on the abstract page https://arxiv.org/abs/2601.03029v2. A full rights notice is delivered at licenses/arxiv-2601.03029-CC-BY-4.0.txt. Characters: every retained excerpt is pure ASCII, so the delivered engine, XeLaTeX with its default Latin Modern text font, reports no missing character.
\begin{lemma}\label{lem:rootofunity}
    Let $K$ be a field. Suppose that $\zeta \in K^\times$ is an element of order~$m < \infty$. Then for all $A, B \in \bar{K}$, we have
    \[
    \prod_{i=1}^m (A - \zeta^iB) = A^m - B^m.
    \]
\end{lemma}

\begin{lemma}\label{lem:div_upgrade}
    Let $\ell$ be prime and $m,n \geq 1$. Let $f \in \mathbb Z[x]$ and assume that $f$ divides $x^n - 1$ modulo $\ell^m$. Then, for any $r \geq 1$, $f$ divides $x^{n \ell^r} - 1$ modulo $\ell^{m+r}$.
\end{lemma}

\begin{lemma}\label{lem:denom_divisible}
     Let $\ell \geq 3$ be prime and assume $\ell \nmid q$.
     Let $s \geq 1$ and $1 \leq t \leq s$. Put $m = m_{\ell,t} = \ell^{t-1}(\ell-1)/2$. Then the polynomial $d_{q,t} := (1+qx^2)^{2m} - x^{2m}$ divides $x^{n_{U}}-1$ modulo $\ell^{s+1-t}$.
\end{lemma}

\begin{proof}
    Let $\alpha \in \mathbb Z$ be a primitive element modulo~$\ell$. Consider the polynomial
    \[
    D_q := \prod_{i=1}^{(\ell-1)/2} (1 + qx^2 - \alpha^i x)(1 + qx^2 + \alpha^i x).
    \]
    Since $\alpha^2$ has order $(\ell-1)/2$ modulo~$\ell$, Lemma~\ref{lem:rootofunity}  yields
    \[
    D_q \equiv (1 + qx^2)^{\ell-1} - x^{\ell-1} \pmod{\ell},
    \]
    and applying Frobenius $t-1$ times gives
    \[
    D_q^{\ell^{t-1}} \equiv d_{q,t} \pmod{\ell}.
    \]
    The polynomial $D_q \pmod{\ell}$ factors as a product of $f_i := 1 - \alpha^i x + qx^2$ for $1 \leq i \leq \ell-1$. Note that $f_i$ and $f_j$ have no roots in common if $i \neq j$.
    
    Assume first that $q$ is a non-square mod~$\ell$. Then none of the $f_i$ is a square. It follows that in this case, $D_q$ is a product of distinct irreducible polynomials of degree at most~2. Since $x^{\ell^2-1}-1 \equiv_{\ell} \prod_{P \in Q} P$, where $Q$ is the set of monic irreducible polynomials of degree $\leq 2$ in $\mathbb F_{\ell}[x]$, we see that $D_q$ divides $x^{\ell^2-1}-1$ mod~$\ell$ and hence $d_{q,t}$ divides $(x^{\ell^2-1}-1)^{\ell^{t-1}} = x^{\ell^{t-1}(\ell^2-1)}-1$ mod~$\ell$.
    We conclude by applying Lemma~\ref{lem:div_upgrade} with $r=s-t$.
    
    Next, assume that $q = r^2$ is a square mod~$\ell$. If we can show that $D_q$ divides $x^{\ell(\ell^2-1)/2}-1$ mod~$\ell$, the result follows in the same way as above. Note that $x^{(\ell^2-1)/2}-1 \equiv_{\ell} \prod_{P \in Q'} P$, where $Q'$ is the set of monic irreducible polynomials of degree $\leq 2$ in $\bF_{\ell}[x]$ whose roots are squares in~$\bF_{\ell^2}$. Thus, we aim to show that the roots of each $f_i$ are squares in~$\bF_{\ell^2}$. If $f_i$ has a root in~$\bF_{\ell}$, this is immediate, so assume $f_i$ is irreducible. It suffices to show that $x$ is a square in $\bF_{\ell^2} = \bF_{\ell}[x]/(f_i)$.

    Solving $x = (Ax + B)^2$ in $\bF_{\ell^2}$ for $A,B$ in $\bF_{\ell}$ yields
    \[
    B^2 - r^{-2}A^2 + x(2AB+r^{-2}\alpha^i A^2-1) = 0.
    \]
    Set $A = \pm rB$. Then the equation has a solution if and only if
    \begin{equation}\label{eq:B}
    B^2(\alpha^i \pm 2r) = 1
    \end{equation}
    has a solution for some $B \in \bF_{\ell}$. Since $f_i$ is irreducible, its discriminant $\alpha^{2i}-4q = (\alpha^i + 2r)(\alpha^i - 2r)$ is a non-square, so precisely one of $\alpha^i + 2r, \alpha^i - 2r$ is a square. Hence \eqref{eq:B} has a solution in $\bF_{\ell}$ and $x$ is a square, as desired.

    Finally, since $q$ is square, two of the $f_i$ are squares of linear factors. Therefore $D_q$ divides $(x^{(\ell^2-1)/2}-1)^{\ell} = x^{\ell(\ell^2-1)/2}-1$ mod~$\ell$, which is what we wanted to show.
\end{proof}

\end{document}
